% !TEX root = ../matan.tex

\section{Линейные векторные пространства. Нормированные пространства, банаховы пространства, примеры}

\subsection*{Линейные векторные пространства}

\begin{Definition}{Линейное (векторное) пространство над $\RR$}{}
	Множество $V$ называется \textbf{линейным пространством} над $\RR$, если на
	нём заданы операции сложения $+:V\times V\to V$ и умножения на скаляр
	$\cdot:\RR\times V\to V$, удовлетворяющие для всех $u,v,w\in V$ и
	$\lambda,\mu\in\RR$ аксиомам:
	\begin{enumerate}
		\item $u+v=v+u$ (коммутативность);
		\item $(u+v)+w=u+(v+w)$ (ассоциативность);
		\item существует нулевой вектор $0\in V$: $v+0=v$;
		\item для каждого $v$ существует противоположный $-v$: $v+(-v)=0$;
		\item $1\cdot v=v$;
		\item $(\lambda\mu)v=\lambda(\mu v)$;
		\item $(\lambda+\mu)v=\lambda v+\mu v$ и $\lambda(u+v)=\lambda u+\lambda v$
		(дистрибутивность).
	\end{enumerate}
\end{Definition}

\begin{Example}{Примеры линейных пространств}{}
	\begin{enumerate}
		\item $\RR^n$ с покоординатными операциями.
		\item Пространство всех многочленов с вещественными коэффициентами.
		\item $C[0,1]$ --- пространство непрерывных функций $f:[0,1]\to\RR$.
		\item Пространство всех вещественных последовательностей
		$x=(x_1,x_2,\dots)$.
	\end{enumerate}
\end{Example}

\subsection*{Нормированные пространства}

\begin{Definition}[required]{Норма}{}
	Пусть $V$ --- линейное пространство над $\RR$. \textbf{Нормой} называется
	отображение $\|\cdot\|:V\to\RR_+$ такое, что для всех $u,v\in V$,
	$\lambda\in\RR$:
	\begin{enumerate}
		\item $\|v\|\ge0$, причём $\|v\|=0\iff v=0$ (невырожденность);
		\item $\|\lambda v\|=|\lambda|\,\|v\|$ (однородность);
		\item $\|u+v\|\le\|u\|+\|v\|$ (неравенство треугольника).
	\end{enumerate}
	Линейное пространство с заданной нормой называется
	\textbf{нормированным пространством}.
\end{Definition}

\begin{Statement}{Норма порождает метрику}{}
	Функция $\rho(x,y):=\|x-y\|$ является метрикой на $V$.
\end{Statement}

\begin{proof}
	\emph{Невырожденность и неотрицательность}: $\rho(x,y)=\|x-y\|\ge0$ и
	$\|x-y\|=0\iff x-y=0\iff x=y$. \emph{Симметрия}:
	$\rho(x,y)=\|x-y\|=|-1|\,\|y-x\|=\|y-x\|=\rho(y,x)$.
	\emph{Неравенство треугольника}: применяя свойство нормы к
	$x-z=(x-y)+(y-z)$,
	\[
	\rho(x,z)=\|x-z\|\le\|x-y\|+\|y-z\|=\rho(x,y)+\rho(y,z). \qedhere
	\]
\end{proof}

\begin{Example}{Нормы в $\RR^n$}{}
	Для $x=(x_1,\dots,x_n)$:
	\[
	\|x\|_p=\left(\sum_{k=1}^{n}|x_k|^p\right)^{1/p}\ (1\le p<\infty),
	\qquad
	\|x\|_\infty=\max_{1\le k\le n}|x_k|.
	\]
	В частности, $\|x\|_2=\sqrt{x_1^2+\dots+x_n^2}$ --- евклидова норма.
	Неравенство треугольника для $\|\cdot\|_p$ --- это неравенство Минковского.
\end{Example}

\begin{Example}{Равномерная норма на $C[0,1]$}{}
	На $C[0,1]$ задаётся \textbf{равномерная} норма
	$\|f\|_\infty=\max_{x\in[0,1]}|f(x)|$. Максимум достигается, так как
	непрерывная функция на отрезке достигает наибольшего значения (теорема
	Вейерштрасса). Сходимость по этой норме --- это равномерная сходимость.
\end{Example}

\subsection*{Банаховы пространства}

\begin{Definition}{Полнота. Банахово пространство}{}
	Нормированное пространство $V$ называется \textbf{полным}, если всякая
	фундаментальная (сходящаяся в себе по Коши) последовательность его
	элементов сходится к элементу того же пространства. Полное нормированное
	пространство называется \textbf{банаховым}.
\end{Definition}

\begin{Example}{Примеры банаховых пространств}{}
	\begin{enumerate}
		\item $\RR^n$ с любой нормой $\|\cdot\|_p$ --- банахово (полнота $\RR^n$).
		\item $C[0,1]$ с равномерной нормой --- банахово: равномерно
		фундаментальная последовательность непрерывных функций равномерно
		сходится к непрерывной функции.
		\item Пространства последовательностей $\ell^p$ (см.\ ниже) при
		$1\le p\le\infty$ --- банаховы.
	\end{enumerate}
\end{Example}

\begin{Example}{Неполное нормированное пространство}{}

	Пространство многочленов на $[0,1]$ с нормой $\|\cdot\|_\infty$ \emph{не}
	полно: последовательность многочленов $p_n(x)=\sum_{k=0}^{n}\frac{x^k}{k!}$
	равномерно фундаментальна, но её предел $e^x$ многочленом не является.
	Таким образом, не всякое нормированное пространство банахово.
\end{Example}

\subsection*{Пространства последовательностей $\ell^p$}

\begin{Definition}{Пространства $\ell^p$ и $\ell^\infty$}{}
	При $1\le p<\infty$
	\[
	\ell^p=\Bigl\{x=(x_1,x_2,\dots)\st \sum_{k=1}^{\infty}|x_k|^p<\infty\Bigr\},
	\qquad
	\|x\|_p=\left(\sum_{k=1}^{\infty}|x_k|^p\right)^{1/p}.
	\]
	Пространство $\ell^\infty$ --- множество ограниченных последовательностей с
	нормой $\|x\|_\infty=\sup_{k}|x_k|$.
\end{Definition}

\begin{Remark}{Супремум вместо максимума}{}
	В $\ell^\infty$ используется именно $\sup$: у ограниченной
	последовательности максимум может не достигаться. Например,
	$x_k=1-\frac1k$ ограничена числом $1$, но ни один член не равен $1$.
\end{Remark}

\subsection*{Ряды в нормированном пространстве}

\begin{Definition}{Сходимость и абсолютная сходимость ряда в $V$}{}
	Пусть $v_k\in V$. Ряд $\sum_{k=1}^{\infty} v_k$ \textbf{сходится} к $S\in V$,
	если $\|S_N-S\|\to0$ для $S_N=\sum_{k=1}^{N} v_k$. Ряд называется
	\textbf{абсолютно сходящимся}, если сходится числовой ряд
	$\sum_{k=1}^{\infty}\|v_k\|$.
\end{Definition}

\begin{Theorem}{Критерий полноты через абсолютную сходимость}{}
	В банаховом пространстве $V$ всякий абсолютно сходящийся ряд сходится.
\end{Theorem}

\begin{proof}
	Пусть $\sum\|v_k\|<\infty$. По критерию Коши для числовых рядов для любого
	$\eps>0$ найдётся $n_0$ такое, что при $n\ge n_0$ и любом $m$ выполнено
	$\sum_{k=n+1}^{n+m}\|v_k\|<\eps$. Тогда для частичных сумм
	\[
	\|S_{n+m}-S_n\|=\left\|\sum_{k=n+1}^{n+m} v_k\right\|
	\le\sum_{k=n+1}^{n+m}\|v_k\|<\eps,
	\]
	то есть $(S_N)$ фундаментальна. В силу полноты $V$ она сходится к
	некоторому $S\in V$, поэтому ряд $\sum v_k$ сходится.
\end{proof}
